\begin{trapbox}
\begin{description}[style=nextline, leftmargin=2.8em, labelsep=0.5em, itemsep=0.35em]
    \item[\textbf{\textcolor{CTrap}{易错点一（误用连续项）}}] 错误认为$m,n,p,q$必须是连续下标，如只敢用$a_1+a_3=a_2+a_2$，忽略$m,n,p,q$可以任意取值。
    \item[\textbf{\textcolor{CTrap}{正确理解（任意匹配）}}] 只要$m+n=p+q$，无论下标是否连续，性质均成立。例如$a_1+a_6 = a_2+a_5 = a_3+a_4$。
    \item[\textbf{\textcolor{CTrap}{错因分析（思维固化）}}] 学生常受“等差中项”影响，只记住相邻项的对称性，忽略了推广到任意下标。
\end{description}

\medskip
\begin{description}[style=nextline, leftmargin=2.8em, labelsep=0.5em, itemsep=0.35em]
    \item[\textbf{\textcolor{CTrap}{易错点二（混淆等比数列）}}] 把等差数列的结论套用到等比数列中，认为等比数列也有$a_m a_n = a_p a_q$（当$m+n=p+q$）。
    \item[\textbf{\textcolor{CTrap}{正确理解（仅等差成立）}}] 该性质只对等差数列成立。等比数列中，下标和相等时，对应项的乘积相等，而非和相等。
    \item[\textbf{\textcolor{CTrap}{错因分析（概念迁移错误）}}] 学生将线性性质与指数性质混淆，未注意等差数列是线性递推，等比数列是指数递推。
\end{description}

\medskip
\begin{description}[style=nextline, leftmargin=2.8em, labelsep=0.5em, itemsep=0.35em]
    \item[\textbf{\textcolor{CTrap}{易错点三（忽略等差数列前提）}}] 认为任意数列都满足$a_m+a_n=a_p+a_q$当$m+n=p+q$。
    \item[\textbf{\textcolor{CTrap}{正确理解（必须等差）}}] 该结论依赖于等差数列的通项线性结构，非等差数列一般不再成立。例如数列$\{n^2\}$中$1^2+3^2=10$，$2^2+2^2=8$，下标和相等但和不相等。
    \item[\textbf{\textcolor{CTrap}{错因分析（记忆泛化）}}] 学生记下结论但忘记前提条件，未检查数列类型便盲目使用。
\end{description}
\end{trapbox}
